\section*{Chapitre 8 \textemdash\ Trigonom\'etrie}
\addcontentsline{toc}{section}{Chapitre 8 \textemdash\ Trigonom\'etrie}

\subsection*{Exercices}
\addcontentsline{toc}{subsection}{Exercices}

%────────────────────────────────────────────────────────────────
\solutiontitle{8}{1}{Conversion degr\'es $\to$ radians}

$\alpha^{\circ}\to\alpha\cdot\pi/180$~:
$30^{\circ}=\frac{\pi}{6}$~; $45^{\circ}=\frac{\pi}{4}$~; $60^{\circ}=\frac{\pi}{3}$~; $90^{\circ}=\frac{\pi}{2}$~;
$120^{\circ}=\frac{2\pi}{3}$~; $135^{\circ}=\frac{3\pi}{4}$~; $150^{\circ}=\frac{5\pi}{6}$~; $180^{\circ}=\pi$~;
$210^{\circ}=\frac{7\pi}{6}$~; $225^{\circ}=\frac{5\pi}{4}$~; $270^{\circ}=\frac{3\pi}{2}$~; $360^{\circ}=2\pi$.

%────────────────────────────────────────────────────────────────
\solutiontitle{8}{2}{Conversion radians $\to$ degr\'es}

$\alpha\,\text{rad}\to\alpha\cdot180/\pi$~:
$\frac{\pi}{3}=60^{\circ}$~; $\frac{\pi}{4}=45^{\circ}$~; $\frac{\pi}{6}=30^{\circ}$~; $\frac{2\pi}{3}=120^{\circ}$~;
$\frac{3\pi}{4}=135^{\circ}$~; $\frac{5\pi}{6}=150^{\circ}$~; $\frac{7\pi}{6}=210^{\circ}$~; $\frac{5\pi}{4}=225^{\circ}$~;
$\frac{4\pi}{3}=240^{\circ}$~; $\frac{3\pi}{2}=270^{\circ}$~; $\frac{7\pi}{4}=315^{\circ}$~; $2\pi=360^{\circ}$.

%────────────────────────────────────────────────────────────────
\solutiontitle{8}{3}{Valeurs exactes}

\begin{enumerate}
  \item $\cos(\frac{\pi}{3})=\frac{1}{2}$.
  \item $\sin(\frac{\pi}{6})=\frac{1}{2}$.
  \item $\tan(\frac{\pi}{4})=1$.
  \item $\cos(\frac{\pi}{2})=0$.
  \item $\sin(\frac{\pi}{2})=1$.
  \item $\cos(\pi)=-1$.
  \item $\sin(\pi)=0$.
  \item $\cos(0)=1$.
  \item $\sin(0)=0$.
  \item $\tan(\frac{\pi}{6})=1/\sqrt{3}=\sqrt{3}/3$.
  \item $\tan(\frac{\pi}{3})=\sqrt{3}$.
  \item $\cos(\frac{\pi}{4})=\frac{\sqrt{2}}{2}$.
\end{enumerate}

%────────────────────────────────────────────────────────────────
\solutiontitle{8}{4}{Quadrant d'un angle}

\begin{enumerate}
  \item $\frac{\pi}{4}\in]0,\frac{\pi}{2}[$~: Q1~; $\cos>0$, $\sin>0$.
  \item $\frac{2\pi}{3}\in]\frac{\pi}{2},\pi[$~: Q2~; $\cos<0$, $\sin>0$.
  \item $\frac{4\pi}{3}\in]\pi,\frac{3\pi}{2}[$~: Q3~; $\cos<0$, $\sin<0$.
  \item $\frac{5\pi}{4}\in]\pi,\frac{3\pi}{2}[$~: Q3~; $\cos<0$, $\sin<0$.
  \item $\frac{7\pi}{6}\in]\pi,\frac{3\pi}{2}[$~: Q3~; $\cos<0$, $\sin<0$.
  \item $\frac{11\pi}{6}\in]\frac{3\pi}{2},2\pi[$~: Q4~; $\cos>0$, $\sin<0$.
\end{enumerate}

%────────────────────────────────────────────────────────────────
\solutiontitle{8}{5}{Sym\'etries du cercle}

\begin{enumerate}
  \item $\cos(\frac{5\pi}{6})=\cos(\pi-\frac{\pi}{6})=-\cos(\frac{\pi}{6})=-\frac{\sqrt{3}}{2}$.
  \item $\sin(\frac{5\pi}{6})=\sin(\pi-\frac{\pi}{6})=\sin(\frac{\pi}{6})=\frac{1}{2}$.
  \item $\cos(\frac{7\pi}{6})=\cos(\pi+\frac{\pi}{6})=-\cos(\frac{\pi}{6})=-\frac{\sqrt{3}}{2}$.
  \item $\sin(\frac{7\pi}{6})=\sin(\pi+\frac{\pi}{6})=-\sin(\frac{\pi}{6})=-\frac{1}{2}$.
  \item $\cos(-\frac{\pi}{3})=\cos(\frac{\pi}{3})=\frac{1}{2}$.
  \item $\sin(-\frac{\pi}{4})=-\sin(\frac{\pi}{4})=-\frac{\sqrt{2}}{2}$.
  \item $\cos(\frac{5\pi}{4})=\cos(\pi+\frac{\pi}{4})=-\cos(\frac{\pi}{4})=-\frac{\sqrt{2}}{2}$.
  \item $\sin(\frac{11\pi}{6})=\sin(2\pi-\frac{\pi}{6})=-\sin(\frac{\pi}{6})=-\frac{1}{2}$.
  \item $\cos(\frac{2\pi}{3})=\cos(\pi-\frac{\pi}{3})=-\cos(\frac{\pi}{3})=-\frac{1}{2}$.
  \item $\sin(\frac{4\pi}{3})=\sin(\pi+\frac{\pi}{3})=-\sin(\frac{\pi}{3})=-\frac{\sqrt{3}}{2}$.
\end{enumerate}

%────────────────────────────────────────────────────────────────
\solutiontitle{8}{6}{P\'eriodicit\'e}

\begin{enumerate}
  \item $\cos(7\frac{\pi}{3})=\cos(7\frac{\pi}{3}-2\pi)=\cos(\frac{\pi}{3})=\frac{1}{2}$.
  \item $\sin(9\frac{\pi}{4})=\sin(9\frac{\pi}{4}-2\pi)=\sin(\frac{\pi}{4})=\frac{\sqrt{2}}{2}$.
  \item $\cos(-\frac{5\pi}{6})=\cos(\frac{5\pi}{6})=-\frac{\sqrt{3}}{2}$.
  \item $\sin(5\pi)=\sin(\pi)=0$.
  \item $\cos(13\frac{\pi}{6})=\cos(13\frac{\pi}{6}-2\pi)=\cos(\frac{\pi}{6})=\frac{\sqrt{3}}{2}$.
  \item $\sin(-\frac{3\pi}{4})=-\sin(\frac{3\pi}{4})=-\frac{\sqrt{2}}{2}$.
  \item $\tan(\frac{5\pi}{4})=\tan(\pi+\frac{\pi}{4})=\tan(\frac{\pi}{4})=1$.
  \item $\cos(100\pi)=\cos(0)=1$ (p\'eriode $2\pi$, $100\pi=50\times2\pi$).
\end{enumerate}

%────────────────────────────────────────────────────────────────
\solutiontitle{8}{7}{Calculer $\cos$ ou $\sin$ \`a partir de l'autre}

\begin{enumerate}
  \item $\sin\theta=3/5$, Q1~: $\cos\theta=\sqrt{1-9/25}=\sqrt{16/25}=4/5$.
  \item $\cos\theta=-\frac{1}{2}$, Q2~: $\sin\theta=\sqrt{1-1/4}=\frac{\sqrt{3}}{2}$.
  \item $\sin\theta=\frac{\sqrt{3}}{2}$, Q2~: $\cos\theta=-\sqrt{1-3/4}=-\frac{1}{2}$.
  \item $\cos\theta=\frac{\sqrt{2}}{2}$, Q4~: $\sin\theta=-\sqrt{1-\frac{1}{2}}=-\frac{\sqrt{2}}{2}$.
  \item $\sin\theta=-1/3$, Q3~: $\cos\theta=-\sqrt{1-1/9}=-\sqrt{8}/3=-2\sqrt{2}/3$.
\end{enumerate}

%────────────────────────────────────────────────────────────────
\solutiontitle{8}{8}{Lire le cercle~: reconstituer l'angle}

\begin{enumerate}
  \item $M(-1,0)$~: $\theta=\pi$.
  \item $M(0,-1)$~: $\theta=\frac{3\pi}{2}$.
  \item $M(\frac{\sqrt{2}}{2},\frac{\sqrt{2}}{2})$~: $\theta=\frac{\pi}{4}$.
  \item $M(\frac{1}{2},-\frac{\sqrt{3}}{2})$~: $\theta=\frac{5\pi}{3}$.
  \item $M(-\frac{\sqrt{3}}{2},\frac{1}{2})$~: $\theta=\frac{5\pi}{6}$.
  \item $M(-\frac{\sqrt{2}}{2},-\frac{\sqrt{2}}{2})$~: $\theta=\frac{5\pi}{4}$.
\end{enumerate}

%────────────────────────────────────────────────────────────────
\solutiontitle{8}{9}{Valeurs de la tangente}

\begin{enumerate}
  \item $\tan(\frac{\pi}{3})=\sqrt{3}$.
  \item $\tan(\frac{2\pi}{3})=\sin(\frac{2\pi}{3})/\cos(\frac{2\pi}{3})=(\frac{\sqrt{3}}{2})/(-\frac{1}{2})=-\sqrt{3}$.
  \item $\tan(\frac{5\pi}{6})=(\frac{1}{2})/(-\frac{\sqrt{3}}{2})=-1/\sqrt{3}=-\sqrt{3}/3$.
  \item $\tan(\frac{5\pi}{4})=(-\frac{\sqrt{2}}{2})/(-\frac{\sqrt{2}}{2})=1$.
  \item $\tan(\frac{7\pi}{6})=(-\frac{1}{2})/(-\frac{\sqrt{3}}{2})=1/\sqrt{3}=\sqrt{3}/3$.
  \item $\tan(\frac{3\pi}{4})=(\frac{\sqrt{2}}{2})/(-\frac{\sqrt{2}}{2})=-1$.
\end{enumerate}

%────────────────────────────────────────────────────────────────
\solutiontitle{8}{10}{Longueur d'arc et angle}

\begin{enumerate}
  \item $\ell=r\theta=5\cdot\frac{\pi}{3}=\frac{5\pi}{3}$~cm.
  \item $\theta=\ell/r=6\pi/9=\frac{2\pi}{3}$~rad~$=120^{\circ}$.
  \item $\ell=6400\cdot\pi/180\approx111{,}7$~km.
  \item $\ell=10\cdot\frac{\pi}{6}=\frac{5\pi}{3}\approx5{,}24$~cm.
\end{enumerate}

%────────────────────────────────────────────────────────────────
\solutiontitle{8}{11}{Lire graphiquement}

\begin{enumerate}
  \item $\cos x=0{,}5=\cos(\frac{\pi}{3})$~: solutions $x=\frac{\pi}{3}$ et $x=2\pi-\frac{\pi}{3}=\frac{5\pi}{3}$.\checkmark
  \item $\sin x=-0{,}5=\sin(-\frac{\pi}{6})$~: sur $[0,2\pi]$~: $x=\pi+\frac{\pi}{6}=\frac{7\pi}{6}$ et $x=2\pi-\frac{\pi}{6}=\frac{11\pi}{6}$.
  \item $\sin x>0$ sur $]0,\pi[$.
  \item $\cos x<0$ sur $]\frac{\pi}{2},\frac{3\pi}{2}[$.
\end{enumerate}

%────────────────────────────────────────────────────────────────
\solutiontitle{8}{12}{Signe de cos et sin selon le quadrant}

\begin{center}
\begin{tabular}{|c|c|c|c|}
\hline
Quadrant & $\cos\theta$ & $\sin\theta$ & $\tan\theta$ \\\hline
Q1 & $+$ & $+$ & $+$ \\
Q2 & $-$ & $+$ & $-$ \\
Q3 & $-$ & $-$ & $+$ \\
Q4 & $+$ & $-$ & $-$ \\\hline
\end{tabular}
\end{center}

%────────────────────────────────────────────────────────────────
\solutiontitle{8}{13}{Cercle trigonom\'etrique et triangle}

Triangle rectangle en $A$~; $\theta=\widehat{ABC}$.
\begin{enumerate}
  \item $\cos\theta=AB/BC$~; $\sin\theta=AC/BC$~; $\tan\theta=AC/AB$.
  \item Triangle $(3,4,5)$, $\theta$ oppos\'e \`a $5$ (i.e. $\theta=90^{\circ}$~?)~:
        Hypot\'enuse $=5$~; c\^ot\'e adjacent \`a $\theta$~? 
        Angle $\theta=\widehat{ABC}$ oppos\'e au c\^ot\'e $AC=?$~.
        Avec $BC=5$ hypot\'enuse~; $AB=4$, $AC=3$~:
        $\cos\theta=4/5$~; $\sin\theta=3/5$~; $\tan\theta=3/4$.
  \item $\cos^2\theta+\sin^2\theta=16/25+9/25=1$.\checkmark
\end{enumerate}

%────────────────────────────────────────────────────────────────
\solutiontitle{8}{14}{R\'esoudre $\cos x=a$ sur $[0,2\pi]$}

\begin{enumerate}
  \item $\cos x=\frac{1}{2}$~: $x=\frac{\pi}{3}$ et $x=\frac{5\pi}{3}$.
  \item $\cos x=-\frac{\sqrt{3}}{2}$~: $x=\frac{5\pi}{6}$ et $x=\frac{7\pi}{6}$.
  \item $\cos x=0$~: $x=\frac{\pi}{2}$ et $x=\frac{3\pi}{2}$.
  \item $\cos x=-1$~: $x=\pi$ (unique).
  \item $\cos x=\frac{\sqrt{2}}{2}$~: $x=\frac{\pi}{4}$ et $x=\frac{7\pi}{4}$.
  \item $\cos x=-\frac{1}{2}$~: $x=\frac{2\pi}{3}$ et $x=\frac{4\pi}{3}$.
  \item $2\cos x-1=0\Rightarrow\cos x=\frac{1}{2}$~: $x=\frac{\pi}{3}$ et $x=\frac{5\pi}{3}$.
  \item $2\cos x+\sqrt{3}=0\Rightarrow\cos x=-\frac{\sqrt{3}}{2}$~: $x=\frac{5\pi}{6}$ et $x=\frac{7\pi}{6}$.
\end{enumerate}

%────────────────────────────────────────────────────────────────
\solutiontitle{8}{15}{R\'esoudre $\sin x=a$ sur $[0,2\pi]$}

\begin{enumerate}
  \item $\sin x=\frac{1}{2}$~: $x=\frac{\pi}{6}$ et $x=\frac{5\pi}{6}$.
  \item $\sin x=-\frac{\sqrt{2}}{2}$~: $x=\frac{5\pi}{4}$ et $x=\frac{7\pi}{4}$.
  \item $\sin x=1$~: $x=\frac{\pi}{2}$ (unique).
  \item $\sin x=0$~: $x=0$ et $x=\pi$ (et $x=2\pi\equiv0$).
  \item $\sin x=\frac{\sqrt{3}}{2}$~: $x=\frac{\pi}{3}$ et $x=\frac{2\pi}{3}$.
  \item $\sin x=-1$~: $x=\frac{3\pi}{2}$ (unique).
  \item $2\sin x-\sqrt{3}=0\Rightarrow\sin x=\frac{\sqrt{3}}{2}$~: $x=\frac{\pi}{3}$ et $x=\frac{2\pi}{3}$.
  \item $\sin x+1=0\Rightarrow\sin x=-1$~: $x=\frac{3\pi}{2}$.
\end{enumerate}

%────────────────────────────────────────────────────────────────
\solutiontitle{8}{16}{R\'esoudre $\tan x=a$ sur $[0,2\pi]$}

\begin{enumerate}
  \item $\tan x=1$~: $x=\frac{\pi}{4}$ et $x=\frac{\pi}{4}+\pi=\frac{5\pi}{4}$.
  \item $\tan x=-1$~: $x=\frac{3\pi}{4}$ et $x=\frac{7\pi}{4}$.
  \item $\tan x=\sqrt{3}$~: $x=\frac{\pi}{3}$ et $x=\frac{4\pi}{3}$.
  \item $\tan x=0$~: $x=0$, $x=\pi$ (et $x=2\pi\equiv0$).
  \item $\tan x=-\sqrt{3}/3=-1/\sqrt{3}$~: $x=\frac{5\pi}{6}$ et $x=\frac{11\pi}{6}$.
  \item $\tan x=\sqrt{3}/3=1/\sqrt{3}$~: $x=\frac{\pi}{6}$ et $x=\frac{7\pi}{6}$.
\end{enumerate}

%────────────────────────────────────────────────────────────────
\solutiontitle{8}{17}{Identit\'es trigonom\'etriques}

\begin{enumerate}
  \item $\frac{1-\cos^2\theta}{\sin\theta}=\frac{\sin^2\theta}{\sin\theta}=\sin\theta$.\hfill$\square$
  \item $\cos^2\theta-\sin^2\theta+2\sin^2\theta=\cos^2\theta+\sin^2\theta=1$.\hfill$\square$
  \item $(\cos\theta+\sin\theta)^2=\cos^2\theta+2\sin\theta\cos\theta+\sin^2\theta=1+2\sin\theta\cos\theta$.\hfill$\square$
  \item $(1-\cos\theta)(1+\cos\theta)=1-\cos^2\theta=\sin^2\theta$.\hfill$\square$
  \item $\cos^4\theta-\sin^4\theta=(\cos^2\theta-\sin^2\theta)(\cos^2\theta+\sin^2\theta)=\cos^2\theta-\sin^2\theta$.\hfill$\square$
\end{enumerate}

%────────────────────────────────────────────────────────────────
\solutiontitle{8}{18}{\'Equations compos\'ees}

\begin{enumerate}
  \item $\cos^2x=1\Rightarrow\cos x=\pm1$~: $x=0,\pi,2\pi$.
  \item $\sin^2x=\frac{1}{2}\Rightarrow\sin x=\pm\frac{\sqrt{2}}{2}$~: $x=\frac{\pi}{4},\frac{3\pi}{4},\frac{5\pi}{4},\frac{7\pi}{4}$.
  \item $\sin x\cdot\cos x=0$~: $\sin x=0$ ou $\cos x=0$~: $x=0,\frac{\pi}{2},\pi,\frac{3\pi}{2},2\pi$.
  \item $2\sin^2x-\sin x=\sin x(2\sin x-1)=0$~: $\sin x=0$ ou $\sin x=\frac{1}{2}$.
        Solutions~: $x=0,\pi,2\pi,\frac{\pi}{6},\frac{5\pi}{6}$.
  \item $\cos^2x-\cos x=\cos x(\cos x-1)=0$~: $\cos x=0$ ou $\cos x=1$.
        Solutions~: $x=0,\frac{\pi}{2},\frac{3\pi}{2},2\pi$.
  \item $(2\cos x-1)(\sin x+1)=0$~: $\cos x=\frac{1}{2}$ ou $\sin x=-1$.
        Solutions~: $x=\frac{\pi}{3},\frac{5\pi}{3},\frac{3\pi}{2}$.
\end{enumerate}

%────────────────────────────────────────────────────────────────
\solutiontitle{8}{19}{Cercle trigonom\'etrique~: lire et placer}

Les $16$ points sont aux angles $0, \frac{\pi}{6}, \frac{\pi}{4}, \ldots, 2\pi$~; voir le cercle du cours.
V\'erification~: $\cos^2(\frac{\pi}{3})+\sin^2(\frac{\pi}{3})=1/4+3/4=1$~;\checkmark
$\cos^2(\frac{\pi}{4})+\sin^2(\frac{\pi}{4})=\frac{1}{2}+\frac{1}{2}=1$~;\checkmark
$\cos^2(\frac{\pi}{6})+\sin^2(\frac{\pi}{6})=3/4+1/4=1$.\checkmark

%────────────────────────────────────────────────────────────────
\solutiontitle{8}{20}{Sym\'etries et valeurs exactes}

\begin{enumerate}
  \item $\cos(\frac{5\pi}{4})+\sin(\frac{7\pi}{6})=-\frac{\sqrt{2}}{2}+(-\frac{1}{2})=-\frac{\sqrt{2}}{2}-\frac{1}{2}$.
  \item $\cos(\frac{2\pi}{3})\times\sin(\frac{5\pi}{3})=(-\frac{1}{2})\times(-\frac{\sqrt{3}}{2})=\sqrt{3}/4$.
  \item $\tan(\frac{4\pi}{3})+\tan(\frac{2\pi}{3})=\sqrt{3}+(-\sqrt{3})=0$.
  \item $\cos^2(\frac{3\pi}{4})+\sin^2(\frac{3\pi}{4})=1$ (relation de Pythagore).
  \item $\sin(\frac{7\pi}{4})-\cos(\frac{5\pi}{6})=-\frac{\sqrt{2}}{2}-(-\frac{\sqrt{3}}{2})=-\frac{\sqrt{2}}{2}+\frac{\sqrt{3}}{2}$.
  \item $\cos(\frac{11\pi}{6})\times\sin(\frac{11\pi}{6})=(\frac{\sqrt{3}}{2})\times(-\frac{1}{2})=-\sqrt{3}/4$.
\end{enumerate}

%────────────────────────────────────────────────────────────────
\solutiontitle{8}{21}{Signe des fonctions et in\'egalit\'es}

\begin{enumerate}
  \item $\sin x>0$ sur $]0,\pi[$.
  \item $\cos x\leq0$ sur $[\frac{\pi}{2},\frac{3\pi}{2}]$.
  \item $\tan x>0$~: Q1 ($]0,\frac{\pi}{2}[$) et Q3 ($]\pi,\frac{3\pi}{2}[$).
  \item $\sin x\geq\cos x\Leftrightarrow\sin x-\cos x\geq0$.
        $\sin x-\cos x=\sqrt{2}\sin(x-\frac{\pi}{4})\geq0\Rightarrow x-\frac{\pi}{4}\in[0,\pi]\Rightarrow x\in[\frac{\pi}{4},\frac{5\pi}{4}]$.
  \item $\sin x+\cos x=\sqrt{2}\sin(x+\frac{\pi}{4})>0\Rightarrow x+\frac{\pi}{4}\in]0,\pi[\Rightarrow x\in]-\frac{\pi}{4},\frac{3\pi}{4}[$~.
        Sur $[0,2\pi]$~: $x\in[0,\frac{3\pi}{4}[\cup]\frac{7\pi}{4},2\pi]$.
\end{enumerate}

%────────────────────────────────────────────────────────────────
\solutiontitle{8}{22}{Calcul de $\tan\theta$ depuis $\sin\theta$ et $\cos\theta$}

\begin{enumerate}
  \item $\sin\theta=4/5$, Q1~: $\cos\theta=3/5$~; $\tan\theta=4/3$.
  \item $\tan\theta=2$, Q1~: $1+\tan^2\theta=5=1/\cos^2\theta\Rightarrow\cos^2\theta=1/5$.
  \item $\cos\theta=-3/5$, Q3~: $\sin\theta=-4/5$~; $\tan\theta=4/3$.
\end{enumerate}

%────────────────────────────────────────────────────────────────
\solutiontitle{8}{23}{Applications physiques}

\begin{enumerate}
  \item $\sin(\frac{\pi}{6})=\frac{1}{2}$~: la composante tangentielle est $mg/2$.
  \item $F_x=100\cos(\frac{\pi}{3})=50$~N~; $F_y=100\sin(\frac{\pi}{3})=50\sqrt{3}\approx86{,}6$~N.
  \item $F_x=F\cos(\frac{\pi}{4})=F/\sqrt{2}=F_y$~: oui, \'egales.\checkmark
\end{enumerate}

%────────────────────────────────────────────────────────────────
\solutiontitle{8}{24}{Tracer $\cos$ et $\sin$}

\begin{enumerate}
  \item Tableau~: voir valeurs remarquables du cours.
  \item Les deux courbes sur $[0,2\pi]$ : sinus et cosinus oscillent entre $-1$ et $1$.
  \item $\cos$~: max $1$ en $0$ (et $2\pi$), min $-1$ en $\pi$, z\'eros en $\frac{\pi}{2}$ et $\frac{3\pi}{2}$.
        $\sin$~: max $1$ en $\frac{\pi}{2}$, min $-1$ en $\frac{3\pi}{2}$, z\'eros en $0$, $\pi$, $2\pi$.
  \item $\cos x\geq\sin x\Leftrightarrow\cos x-\sin x\geq0$~: $x\in[0,\frac{\pi}{4}]\cup[\frac{5\pi}{4},2\pi]$.
\end{enumerate}

%────────────────────────────────────────────────────────────────
\solutiontitle{8}{25}{R\'esoudre une \'equation par substitution}

\begin{enumerate}
  \item $u=\sin x$~: $2u^2+u-1=(2u-1)(u+1)=0\Rightarrow u=\frac{1}{2}$ ou $u=-1$.
        $\sin x=\frac{1}{2}$~: $x=\frac{\pi}{6},\frac{5\pi}{6}$~; $\sin x=-1$~: $x=\frac{3\pi}{2}$.
  \item $u=\cos x$~: $2u^2-u=u(2u-1)=0\Rightarrow u=0$ ou $u=\frac{1}{2}$.
        $\cos x=0$~: $x=\frac{\pi}{2},\frac{3\pi}{2}$~; $\cos x=\frac{1}{2}$~: $x=\frac{\pi}{3},\frac{5\pi}{3}$.
  \item $4\sin^2x-3=0\Rightarrow\sin^2x=3/4\Rightarrow\sin x=\pm\frac{\sqrt{3}}{2}$.
        Solutions~: $x=\frac{\pi}{3},\frac{2\pi}{3},\frac{4\pi}{3},\frac{5\pi}{3}$.
  \item $u=\cos x$~: $u^2+2u+1=(u+1)^2=0\Rightarrow u=-1\Rightarrow\cos x=-1$~: $x=\pi$.
\end{enumerate}

%────────────────────────────────────────────────────────────────
\solutiontitle{8}{26}{Signe et quadrant combin\'es}

\begin{enumerate}
  \item $\cos>0$ et $\sin<0$~: Q4~: $\theta\in]\frac{3\pi}{2},2\pi[$.
  \item $\cos<0$ et $\sin<0$~: Q3~: $\theta\in]\pi,\frac{3\pi}{2}[$.
  \item $\tan>0$ et $\cos<0$~: Q3 (tan positif en Q1 et Q3, cos n\'egatif en Q2 et Q3)~: Q3.
  \item $\cos\theta=\frac{\sqrt{3}}{2}$~: $\theta=\frac{\pi}{6}$ ou $\theta=\frac{11\pi}{6}$. Avec $\sin\theta>0$~: $\theta=\frac{\pi}{6}$.
  \item $\sin\theta=-\frac{1}{2}$~: $\theta=\frac{7\pi}{6}$ ou $\theta=\frac{11\pi}{6}$. Avec $\cos\theta<0$~: $\theta=\frac{7\pi}{6}$.
\end{enumerate}

%────────────────────────────────────────────────────────────────
\solutiontitle{8}{27}{\'Equations m\^el\'ees}

\begin{enumerate}
  \item $\sin x=\cos x\Rightarrow\tan x=1\Rightarrow x=\frac{\pi}{4}$ ou $x=\frac{5\pi}{4}$.
  \item $\sin x+\cos x=0\Rightarrow\tan x=-1\Rightarrow x=\frac{3\pi}{4}$ ou $x=\frac{7\pi}{4}$.
  \item $\sin^2x=\cos^2x\Rightarrow\tan^2x=1\Rightarrow\tan x=\pm1$~: $x=\frac{\pi}{4},\frac{3\pi}{4},\frac{5\pi}{4},\frac{7\pi}{4}$.
  \item $2\sin x\cos x=1=\sin(2x)$ (adm.)~: $2x=\frac{\pi}{2}$ ou $2x=5\frac{\pi}{2}$~: $x=\frac{\pi}{4}$ ou $x=\frac{5\pi}{4}$.
\end{enumerate}

%────────────────────────────────────────────────────────────────
\solutiontitle{8}{28}{Valeurs num\'eriques exactes}

\begin{enumerate}
  \item $\cos^2(\frac{\pi}{6})+\sin^2(\frac{\pi}{6})=3/4+1/4=1$.
  \item $\cos(\frac{2\pi}{3})+2\sin(\frac{\pi}{6})=-\frac{1}{2}+2\cdot\frac{1}{2}=0$.
  \item $\tan(\frac{\pi}{4})\cdot\cos(\frac{\pi}{4})=1\cdot\frac{\sqrt{2}}{2}=\frac{\sqrt{2}}{2}$.
  \item $2\cos^2(\frac{\pi}{3})-1=2\cdot1/4-1=-\frac{1}{2}$.
  \item $\sin(\frac{\pi}{4})\cos(\frac{\pi}{4})=(\frac{\sqrt{2}}{2})^2=\frac{1}{2}$.
  \item $\tan(\frac{\pi}{3})\cdot\sin(\frac{\pi}{6})=\sqrt{3}\cdot\frac{1}{2}=\frac{\sqrt{3}}{2}$.
\end{enumerate}

%────────────────────────────────────────────────────────────────
\solutiontitle{8}{29}{Trigonom\'etrie dans un triangle}

$\widehat{A}=\frac{\pi}{3}$, $\widehat{B}=\frac{\pi}{4}$, $AB=6$.
\begin{enumerate}
  \item $\widehat{C}=\pi-\frac{\pi}{3}-\frac{\pi}{4}=\pi-7\pi/12=5\pi/12=75^{\circ}$.
  \item $\widehat{A}+\widehat{B}+\widehat{C}=\frac{\pi}{3}+\frac{\pi}{4}+5\pi/12=4\pi/12+3\pi/12+5\pi/12=12\pi/12=\pi$.\checkmark
        Non rectangle (pas d'angle $\frac{\pi}{2}$).
  \item Si rectangle en $C$ (hypothétiquement)~: $\widehat{C}=\frac{\pi}{2}\Rightarrow\widehat{A}+\widehat{B}=\frac{\pi}{2}$~;
        mais $\frac{\pi}{3}+\frac{\pi}{4}=7\pi/12\ne\frac{\pi}{2}$~: incompatible avec les angles donn\'es.
\end{enumerate}

%────────────────────────────────────────────────────────────────
\solutiontitle{8}{30}{Fonctions trigonom\'etriques~: variations}

\begin{enumerate}
  \item $\cos\theta$ croissante sur $[\pi,2\pi]$, d\'ecroissante sur $[0,\pi]$.
  \item $\sin\theta$ d\'ecroissante sur $[\frac{\pi}{2},\frac{3\pi}{2}]$.
  \item $\cos$~: max $1$ en $0$ et $2\pi$~; min $-1$ en $\pi$.
        $\sin$~: max $1$ en $\frac{\pi}{2}$~; min $-1$ en $\frac{3\pi}{2}$.
  \item Ant\'ec\'edents de $\frac{1}{2}$ par $\cos$~: $\frac{\pi}{3}$ et $\frac{5\pi}{3}$.
        Par $\sin$~: $\frac{\pi}{6}$ et $\frac{5\pi}{6}$.
\end{enumerate}

%────────────────────────────────────────────────────────────────
\solutiontitle{8}{31}{Calculer une expression}

\begin{enumerate}
  \item $\sin^2\theta+\cos^2\theta+\tan^2\theta\cos^2\theta=1+\sin^2\theta=1+\sin^2\theta$.
        Plus pr\'ecis\'ement~: $\tan^2\theta\cos^2\theta=\sin^2\theta$, donc total~$=1+\sin^2\theta$.
        Hmm~: $\sin^2\theta+\cos^2\theta+\tan^2\theta\cos^2\theta=1+\frac{\sin^2\theta}{\cos^2\theta}\cdot\cos^2\theta=1+\sin^2\theta$.
  \item $(\sin\theta-1)(\sin\theta+1)=\sin^2\theta-1=-\cos^2\theta$.
  \item $\frac{\sin^2\theta-1}{\cos\theta}=\frac{-\cos^2\theta}{\cos\theta}=-\cos\theta$.
  \item $(\cos+\sin)^2+(\cos-\sin)^2=2\cos^2\theta+2\sin^2\theta=2$.
  \item $\frac{1}{\cos^2\theta}-\tan^2\theta=\frac{1}{\cos^2\theta}-\frac{\sin^2\theta}{\cos^2\theta}=\frac{1-\sin^2\theta}{\cos^2\theta}=\frac{\cos^2\theta}{\cos^2\theta}=1$.
\end{enumerate}

%────────────────────────────────────────────────────────────────
\solutiontitle{8}{32}{R\'esoudre avec les sym\'etries}

\begin{enumerate}
  \item $\cos(\pi-x)=-\cos x=-\frac{1}{2}\Rightarrow\cos x=\frac{1}{2}\Rightarrow x=\frac{\pi}{3}$ ou $x=\frac{5\pi}{3}$.
  \item $\sin(\pi-x)=\sin x=\frac{\sqrt{3}}{2}\Rightarrow x=\frac{\pi}{3}$ ou $x=\frac{2\pi}{3}$.
  \item $\cos(-x)=\cos x=\frac{\sqrt{2}}{2}\Rightarrow x=\frac{\pi}{4}$ ou $x=\frac{7\pi}{4}$.
  \item $\sin(\pi+x)=-\sin x=\frac{1}{2}\Rightarrow\sin x=-\frac{1}{2}\Rightarrow x=\frac{7\pi}{6}$ ou $x=\frac{11\pi}{6}$.
  \item $\cos(x+\pi)=-\cos x=-\frac{\sqrt{3}}{2}\Rightarrow\cos x=\frac{\sqrt{3}}{2}\Rightarrow x=\frac{\pi}{6}$ ou $x=\frac{11\pi}{6}$.
\end{enumerate}

%────────────────────────────────────────────────────────────────
\solutiontitle{8}{33}{\logicoalgo\ Algorithme~: r\'esoudre $\cos x=a$ sur $[0,2\pi]$}

\begin{enumerate}
  \item Si $\alpha=0$~: les deux solutions $\alpha$ et $2\pi-\alpha=2\pi\equiv0$ co\"incident.
        Si $\alpha=\pi$~: $2\pi-\alpha=\pi$~: idem. Logiquement~: $\alpha=2\pi-\alpha\Leftrightarrow\alpha=k\pi$.\hfill$\square$
  \item $a=\frac{1}{2}$~: $\alpha=\frac{\pi}{3}$~; deux solutions $\frac{\pi}{3}$ et $\frac{5\pi}{3}$.
        $a=-1$~: $\alpha=\pi$~; unique solution $\pi$.
        $a=\frac{\sqrt{2}}{2}$~: $\alpha=\frac{\pi}{4}$~; deux solutions $\frac{\pi}{4}$ et $\frac{7\pi}{4}$.
  \item Analogue pour $\sin x=a$~: $\alpha=\arcsin(a)\in[-\frac{\pi}{2},\frac{\pi}{2}]$~;
        adapter \`a $[0,2\pi]$ si $\alpha<0$ (ajouter $2\pi$)~;
        deuxi\`eme solution~: $\pi-\alpha$ (adap. si n\'ecessaire). Diff\'erence~: angles de r\'ef\'erence.
  \item
\begin{lstlisting}[language=Python]
import math
def solve_cos(a):
    if abs(a) > 1:
        print("Pas de solution")
        return
    alpha = math.acos(a)
    if alpha == 0 or abs(alpha - math.pi) < 1e-12:
        print(f"Solution unique: x = {alpha:.6f}")
    else:
        print(f"x1={alpha:.6f}, x2={2*math.pi-alpha:.6f}")
\end{lstlisting}
\end{enumerate}

%────────────────────────────────────────────────────────────────
\solutiontitle{8}{34}{Valeurs exactes par le cercle}

\begin{enumerate}
  \item $\cos(\frac{3\pi}{4})=-\cos(\frac{\pi}{4})=-\frac{\sqrt{2}}{2}$~;
        $\sin(\frac{5\pi}{6})=\sin(\frac{\pi}{6})=\frac{1}{2}$~;
        $\cos(\frac{7\pi}{6})=-\cos(\frac{\pi}{6})=-\frac{\sqrt{3}}{2}$~;
        $\sin(\frac{11\pi}{6})=-\sin(\frac{\pi}{6})=-\frac{1}{2}$.
  \item $\cos^2(\frac{5\pi}{4})+\sin^2(\frac{5\pi}{4})=(-\frac{\sqrt{2}}{2})^2+(-\frac{\sqrt{2}}{2})^2=\frac{1}{2}+\frac{1}{2}=1$.\checkmark
  \item $\cos^2\theta=1-3/4=1/4$~; $\theta\in]\frac{\pi}{2},\pi[$~: $\cos\theta<0$~: $\cos\theta=-\frac{1}{2}$.
  \item $\cos\theta=\sin\theta\Rightarrow\tan\theta=1\Rightarrow\theta=\frac{\pi}{4}$ ou $\theta=\frac{5\pi}{4}$.
\end{enumerate}

%────────────────────────────────────────────────────────────────
\solutiontitle{8}{35}{Cercle trigonom\'etrique et g\'eom\'etrie}

\begin{enumerate}
  \item $MN^2=(\cos\theta-\cos\varphi)^2+(\sin\theta-\sin\varphi)^2
        =\cos^2\theta-2\cos\theta\cos\varphi+\cos^2\varphi+\sin^2\theta-2\sin\theta\sin\varphi+\sin^2\varphi
        =2-2(\cos\theta\cos\varphi+\sin\theta\sin\varphi)$.
  \item $\theta=\frac{\pi}{3}$, $\varphi=0$~: $MN^2=2-2\cos(\frac{\pi}{3})=2-1=1$~; $MN=1$.
  \item $\theta=\frac{\pi}{2}$, $\varphi=\frac{\pi}{6}$~: $MN^2=2-2(0\cdot\cos(\frac{\pi}{6})+1\cdot\sin(\frac{\pi}{6}))=2-2\cdot\frac{1}{2}=1$~; $MN=1$.
  \item Diam\`etre~: $\theta=0$, $\varphi=\pi$~: $MN^2=2-2(-1)=4$~; $MN=2$.\checkmark (rayon $1$, diam\'etre $2$).
\end{enumerate}

%────────────────────────────────────────────────────────────────
\solutiontitle{8}{36}{Signe et quadrant~: discussion}

\begin{enumerate}
  \item $\cos\theta<0$ et $\sin\theta>0$~: Q2.
  \item $\theta\in]\frac{\pi}{2},\pi[$.
  \item $\tan\theta=-1$ et Q2~: $\theta=\frac{3\pi}{4}$.
  \item $\cos\theta=-\frac{\sqrt{3}}{2}$~: $\theta=\frac{5\pi}{6}$ ou $\theta=\frac{7\pi}{6}$~;
        avec $\sin\theta>0$~: $\theta=\frac{5\pi}{6}$ (Q2).
\end{enumerate}

%────────────────────────────────────────────────────────────────
\solutiontitle{8}{37}{Identit\'es et factorisation}

\begin{enumerate}
  \item $\cos^2\theta-\sin^2\theta=(\cos\theta-\sin\theta)(\cos\theta+\sin\theta)$.
        $2\cos^2\theta-1=\cos^2\theta-\sin^2\theta$ (car $\sin^2\theta=1-\cos^2\theta$).\hfill$\square$
  \item $\frac{\sin\theta}{\cos\theta}+\frac{\cos\theta}{\sin\theta}
        =\frac{\sin^2\theta+\cos^2\theta}{\sin\theta\cos\theta}=\frac{1}{\sin\theta\cos\theta}$.\hfill$\square$
  \item $(\cos\theta-\sin\theta)^2+(\cos\theta+\sin\theta)^2
        =2\cos^2\theta+2\sin^2\theta=2$.\hfill$\square$
  \item $1-2\sin^2\theta=1-2(1-\cos^2\theta)=2\cos^2\theta-1$.\hfill$\square$
\end{enumerate}

%────────────────────────────────────────────────────────────────
\solutiontitle{8}{38}{\logiq\ Propositions sur $\cos$ et $\sin$}

\begin{enumerate}
  \item \textbf{Vraie.} N\'egation~: $\exists\theta\in\R,\;\cos^2\theta+\sin^2\theta\ne1$~: fausse.
  \item \textbf{Fausse.} $\cos\theta\in[-1,1]$ pour tout $\theta$. N\'egation~:
        $\forall\theta,\;\cos\theta\leq1$~: vraie.
  \item \textbf{Vraie.} $\sin\theta\geq-1$ toujours. N\'egation~: $\exists\theta,\;\sin\theta<-1$~: fausse.
  \item \textbf{Fausse.} Sur $[0,\pi]$, $\cos\theta\in[-1,1]$, donc $\cos\theta=2$ est impossible.
        N\'egation~: $\forall\theta\in[0,\pi],\;\cos\theta\ne2$~: vraie.
\end{enumerate}

%────────────────────────────────────────────────────────────────
\solutiontitle{8}{39}{\logiq\ \'Equivalences trigonom\'etriques}

\begin{enumerate}
  \item $\Leftrightarrow$ (d\'efinition des z\'eros du cosinus).
  \item $\Leftarrow$~: $\theta\in]0,\pi[\Rightarrow\sin\theta>0$~: vraie.
        $\Rightarrow$~: sur $[0,2\pi]$, $\sin\theta>0$ implique $\theta\in]0,\pi[$~: vraie.
        Donc $\Leftrightarrow$ sur $[0,2\pi]$.
  \item $\Leftrightarrow$ (d\'efinition de l'\'egalit\'e des cosinus).
  \item $\Leftrightarrow$ (sur $\R$ via la p\'eriodicit\'e).
\end{enumerate}

%────────────────────────────────────────────────────────────────
\solutiontitle{8}{40}{\logiq\ Conditions sur les solutions d'une \'equation trig.}

\begin{enumerate}
  \item $\cos x=a$ a des solutions ssi $|a|\leq1$.
        Formuler~: $\exists x\in\R,\;\cos x=a\Leftrightarrow a\in[-1,1]$.
  \item Solution unique dans $[0,2\pi]$~: $a=1$ ($x=0\equiv2\pi$) ou $a=-1$ ($x=\pi$).
        Condition~: $a=1$ ou $a=-1$.
  \item N\'egation de \og deux solutions\fg~:
        l'\'equation a $0$ ou $1$ solution, i.e. $|a|>1$ ou $a=\pm1$.
\end{enumerate}

%────────────────────────────────────────────────────────────────
\solutiontitle{8}{41}{\logiq\ Raisonnement sur les sym\'etries du cercle}

\begin{enumerate}
  \item Sur le cercle~: $-\theta$ et $\theta$ donnent des points sym\'etriques par rapport \`a $Ox$~:
        m\^eme abscisse, donc $\cos(-\theta)=\cos\theta$. Raisonnement \textbf{g\'eom\'etrique direct}.\hfill$\square$
  \item Le point associ\'e \`a $\pi-\theta$ est le sym\'etrique de celui associ\'e \`a $\theta$
        par rapport \`a $Oy$~: m\^eme ordonn\'ee, donc $\sin(\pi-\theta)=\sin\theta$.\hfill$\square$
  \item Implication~: \og Pour tout $\theta$, $\cos(\pi+\theta)=-\cos\theta$.\fg
        R\'eciproque~: \og Si $\cos(\pi+\theta)=-\cos\theta$, alors \ldots\fg~: toujours vraie (identit\'e pour tout $\theta$).
  \item $(\cos\theta+\sin\theta)^2=1+2\sin\theta\cos\theta\leq1+2\cdot\frac{1}{2}=2$
        (car $2\sin\theta\cos\theta=\sin(2\theta)\in[-1,1]$).\hfill$\square$
\end{enumerate}

%────────────────────────────────────────────────────────────────
\solutiontitle{8}{42}{\logiq\ Propositions combinant quadrant et valeurs}

\begin{enumerate}
  \item Implication vraie. R\'eciproque~: si $\cos\theta<0$ et $\sin\theta>0$, est-ce que $\theta\in]\frac{\pi}{2},\pi[$~?
        Sur $[0,2\pi]$, oui~: c'est la d\'efinition du Q2. \textbf{\'Equivalence} sur $[0,2\pi]$.
  \item $\cos\theta<0$ et $\sin\theta>0\Leftrightarrow\theta\in]\frac{\pi}{2},\pi[$.
  \item $\tan\theta>0\Leftrightarrow\sin\theta/\cos\theta>0\Leftrightarrow(\sin\theta>0\wedge\cos\theta>0)\vee(\sin\theta<0\wedge\cos\theta<0)$.\hfill$\square$
\end{enumerate}

%────────────────────────────────────────────────────────────────
\solutiontitle{8}{43}{\logiq\ N\'ecessaire et suffisant en trigonom\'etrie}

\begin{enumerate}
  \item $\sin\theta=0$ est \textbf{CN} pour $\theta=0$ (car $\sin 0=0$, OK~; mais $\sin\pi=0\ne$ implique $\theta=0$).
        CNS pour $\theta\in\{0,\pi,2\pi\}$~? Non, $\sin(\theta)=0\Leftrightarrow\theta\in\{0,\pi,2\pi,\ldots\}$.
        CNS pour $\theta=k\pi$~: oui.
  \item $\sin\theta=\frac{\sqrt{3}}{2}$ est \textbf{CS} (pas CNS) pour $\theta=\frac{\pi}{3}$~:
        l'implication $\theta=\frac{\pi}{3}\Rightarrow\sin\theta=\frac{\sqrt{3}}{2}$ est vraie~;
        mais $\sin(\frac{2\pi}{3})=\frac{\sqrt{3}}{2}$ aussi, donc la CNS est $\theta=\frac{\pi}{3}$ ou $\theta=\frac{2\pi}{3}$.
  \item CNS pour $\tan\theta=1$ sur $[0,2\pi]$~: $\theta=\frac{\pi}{4}$ ou $\theta=\frac{5\pi}{4}$.
\end{enumerate}

%────────────────────────────────────────────────────────────────
\solutiontitle{8}{44}{\algo\ Algorithme~: table de valeurs de $\cos$ et $\sin$ \computer}

\begin{lstlisting}[language=Python]
import math
angles = [0, math.pi/6, math.pi/4, math.pi/3, math.pi/2,
          2*math.pi/3, 3*math.pi/4, 5*math.pi/6, math.pi,
          7*math.pi/6, 5*math.pi/4, 4*math.pi/3, 3*math.pi/2,
          5*math.pi/3, 7*math.pi/4, 11*math.pi/6, 2*math.pi]
for t in angles:
    c, s = math.cos(t), math.sin(t)
    verif = round(c**2 + s**2, 10)
    print(f"theta={t:.4f}: cos={c:.4f}, sin={s:.4f}, sum={verif}")
\end{lstlisting}

%────────────────────────────────────────────────────────────────
\solutiontitle{8}{45}{\algo\ Algorithme~: r\'esoudre $\cos x=a$ sur $[0,2\pi]$}

\begin{lstlisting}[language=Python]
import math

def solve_cos_sin(a, mode='cos'):
    if abs(a) > 1:
        print("Pas de solution")
        return
    if mode == 'cos':
        alpha = math.acos(a)
        s1, s2 = alpha, 2*math.pi - alpha
    else:
        alpha = math.asin(a)
        if alpha < 0: alpha += 2*math.pi
        s1, s2 = alpha, math.pi - alpha
        if s2 < 0: s2 += 2*math.pi
    print(f"x1 = {s1:.6f} rad = {math.degrees(s1):.2f} deg")
    if abs(s1 - s2) > 1e-10:
        print(f"x2 = {s2:.6f} rad = {math.degrees(s2):.2f} deg")

solve_cos_sin(0.5, 'cos')
solve_cos_sin(-1, 'cos')
solve_cos_sin(0.5, 'sin')
\end{lstlisting}

%────────────────────────────────────────────────────────────────
\solutiontitle{8}{46}{\algo\ Algorithme~: tracer le cercle trigonom\'etrique \computer}

\begin{lstlisting}[language=Python]
import math, matplotlib.pyplot as plt
import numpy as np

fig, ax = plt.subplots(figsize=(8,8))
theta = np.linspace(0, 2*math.pi, 300)
ax.plot(np.cos(theta), np.sin(theta), 'b-')
ax.set_aspect('equal')
ax.axhline(0, color='k', lw=0.5); ax.axvline(0, color='k', lw=0.5)
# 12 angles remarquables
angles = [k*math.pi/6 for k in range(12)]
for t in angles:
    x, y = math.cos(t), math.sin(t)
    ax.plot(x, y, 'ro', ms=5)
    ax.annotate(f'{t:.2f}', (x*1.15, y*1.15), ha='center', fontsize=8)
# Point M pour theta saisi
t0 = math.pi/3
ax.plot(math.cos(t0), math.sin(t0), 'g*', ms=15)
ax.plot([0, math.cos(t0)], [0, math.sin(t0)], 'g--')
ax.plot([math.cos(t0), math.cos(t0)], [0, math.sin(t0)], 'r--')
ax.plot([0, math.cos(t0)], [math.sin(t0), math.sin(t0)], 'b--')
plt.title('Cercle trigonometrique'); plt.show()
\end{lstlisting}

%────────────────────────────────────────────────────────────────
\solutiontitle{8}{47}{\algo\ Algorithme~: \'equations trigonom\'etriques par balayage \computer}

\begin{lstlisting}[language=Python]
import math

f = lambda x: 2*math.sin(x)**2 - 3*math.sin(x) + 1
a, b, h = 0, 2*math.pi, 0.001
prev, x = f(a), a
while x + h <= b:
    curr = f(x + h)
    if prev * curr < 0:
        print(f"Racine dans [{x:.3f}, {x+h:.3f}]")
    prev, x = curr, x + h
# Solutions exactes: sin x = \frac{1}{2} (x=pi/6, 5pi/6) ou sin x = 1 (x=pi/2)
\end{lstlisting}

%────────────────────────────────────────────────────────────────
\solutiontitle{8}{48}{\algo\ Algorithme~: longueur de la corde de l'angle $\theta$ \computer}

\begin{lstlisting}[language=Python]
import math

def corde(theta):
    # Distance entre (cos(0),sin(0))=(1,0) et (cos(t),sin(t))
    return math.sqrt((math.cos(theta)-1)**2 + math.sin(theta)**2)

for n in [2, 4, 6, 8, 12, 100]:
    t = math.pi / n
    print(f"n={n:4d}: theta={t:.4f}, corde={corde(t):.8f}")
# La corde tend vers 0 quand theta -> 0
\end{lstlisting}

%────────────────────────────────────────────────────────────────
\solutiontitle{8}{49}{\algo\ Algorithme~: nombre $\pi$ par p\'erim\`etre de polygone \computer}

\begin{lstlisting}[language=Python]
import math

for n in [4, 6, 8, 12, 24, 96, 1000]:
    Pn = 2 * n * math.sin(math.pi / n)
    print(f"n={n:5d}: P_n = {Pn:.8f}, pi = {math.pi:.8f}")
# Pour n=96: P_96 ~ 3.14103... (approx d'Archimede)
\end{lstlisting}

%────────────────────────────────────────────────────────────────
\subsection*{Probl\`emes}
\addcontentsline{toc}{subsection}{Probl\`emes}

%────────────────────────────────────────────────────────────────
\psolutiontitle{8}{1}{Trigonom\'etrie dans un triangle rectangle}

Triangle $ABC$ rectangle en $A$, $BC=10$, $\widehat{ABC}=\frac{\pi}{6}$.
\textbf{(1)} $AB=BC\cos(\widehat{ABC})=10\cos(\frac{\pi}{6})=10\cdot\frac{\sqrt{3}}{2}=5\sqrt{3}$.
$AC=BC\sin(\widehat{ABC})=10\sin(\frac{\pi}{6})=5$.
\textbf{(2)} $AB^2+AC^2=75+25=100=BC^2$.\checkmark
\textbf{(3)} $\tan(\widehat{ABC})=AC/AB=5/(5\sqrt{3})=1/\sqrt{3}=\tan(\frac{\pi}{6})$.\checkmark
$\tan(\widehat{ACB})=AB/AC=\sqrt{3}=\tan(\frac{\pi}{3})$.
\textbf{(4)} $\widehat{ACB}=\pi-\frac{\pi}{2}-\frac{\pi}{6}=\frac{\pi}{3}$.\checkmark

%────────────────────────────────────────────────────────────────
\psolutiontitle{8}{2}{R\'esolution syst\'ematique}

\textbf{(1)} $u=\sin x$~: $2u^2-3u+1=(2u-1)(u-1)=0\Rightarrow u=\frac{1}{2}$ ou $u=1$.
$\sin x=\frac{1}{2}$~: $x=\frac{\pi}{6},\frac{5\pi}{6}$~; $\sin x=1$~: $x=\frac{\pi}{2}$.

\textbf{(2)} $\cos^2x-\sin^2x=0\Rightarrow(\cos x-\sin x)(\cos x+\sin x)=0$.
$\cos x=\sin x\Rightarrow\tan x=1$~: $x=\frac{\pi}{4},\frac{5\pi}{4}$~;
$\cos x=-\sin x\Rightarrow\tan x=-1$~: $x=\frac{3\pi}{4},\frac{7\pi}{4}$.

\textbf{(3)} $\sin x=\cos x+1$~: \'elever au carr\'e~: $\sin^2x=\cos^2x+2\cos x+1$~;
$(1-\cos^2x)=\cos^2x+2\cos x+1$~; $2\cos^2x+2\cos x=0$~; $\cos x(\cos x+1)=0$.
$\cos x=0$~: $x=\frac{\pi}{2},\frac{3\pi}{2}$~; $\cos x=-1$~: $x=\pi$.
V\'erif.~: $x=\frac{\pi}{2}$~: $\sin(\frac{\pi}{2})=1$, $\cos(\frac{\pi}{2})+1=1$.\checkmark
$x=\frac{3\pi}{2}$~: $\sin=-1$, $\cos+1=1\ne-1$~: solution parasite.\quad
$x=\pi$~: $\sin=0$, $\cos+1=0$.\checkmark
Solutions~: $x=\frac{\pi}{2}$ et $x=\pi$.

\textbf{(4)} $u=\cos x$~: $2u^2+u-1=(2u-1)(u+1)=0\Rightarrow u=\frac{1}{2}$ ou $u=-1$.
$\cos x=\frac{1}{2}$~: $x=\frac{\pi}{3},\frac{5\pi}{3}$~; $\cos x=-1$~: $x=\pi$.

%────────────────────────────────────────────────────────────────
\psolutiontitle{8}{3}{Mesure d'une hauteur inaccessible}

Angle $\alpha=\frac{\pi}{4}$ \`a $20$~m, $\beta=\frac{\pi}{3}$ \`a $10$~m du pied de l'arbre.
\textbf{(1)} Soit $d$ la distance horiz. du point B au pied~:
$\tan(\frac{\pi}{4})=h/(d+10)=1\Rightarrow h=d+10$~;
$\tan(\frac{\pi}{3})=h/d=\sqrt{3}\Rightarrow h=d\sqrt{3}$.
\textbf{(2)} $d+10=d\sqrt{3}\Rightarrow d(\sqrt{3}-1)=10\Rightarrow d=\frac{10}{\sqrt{3}-1}=\frac{10(\sqrt{3}+1)}{2}=5(\sqrt{3}+1)$.
$h=5\sqrt{3}(\sqrt{3}+1)=5(3+\sqrt{3})=15+5\sqrt{3}$.
\textbf{(3)} $h=15+5\sqrt{3}\approx15+8{,}66=23{,}66$~m.

%────────────────────────────────────────────────────────────────
\psolutiontitle{8}{4}{Coordonn\'ees sur le cercle}

$M(\cos\theta,\sin\theta)$.
\textbf{(1)} $MA^2=(\cos\theta-1)^2+\sin^2\theta=\cos^2\theta-2\cos\theta+1+\sin^2\theta=2-2\cos\theta$.
$MA=\sqrt{2-2\cos\theta}$.
\textbf{(2)} $MA=1\Rightarrow2-2\cos\theta=1\Rightarrow\cos\theta=\frac{1}{2}\Rightarrow\theta=\frac{\pi}{3}$ (sur $[0,\pi]$).
\textbf{(3)} $MB^2=\cos^2\theta+(\sin\theta-1)^2=2-2\sin\theta$~; $MB=\sqrt{2-2\sin\theta}$.
\textbf{(4)} $MB=1\Rightarrow\sin\theta=\frac{1}{2}\Rightarrow\theta=\frac{\pi}{6}$ (sur $[0,\pi]$).

%────────────────────────────────────────────────────────────────
\psolutiontitle{8}{5}{Applications des formules}

\textbf{(1)} $\sin\theta=0{,}6$, Q1~:
$\cos\theta=0{,}8$~; $\tan\theta=0{,}75$~;
$\sin(\pi-\theta)=\sin\theta=0{,}6$~;
$\cos(\pi+\theta)=-\cos\theta=-0{,}8$.

\textbf{(2)} $\cos\theta=-0{,}8$, Q3~:
$\sin\theta=-\sqrt{1-0{,}64}=-0{,}6$~; $\tan\theta=0{,}75$~;
$\cos(-\theta)=\cos\theta=-0{,}8$~;
$\sin(\pi-\theta)=\sin\theta=-0{,}6$.

%────────────────────────────────────────────────────────────────
\psolutiontitle{8}{6}{Trigonom\'etrie et coordonn\'ees}

$\cos\theta=3/5$.
\textbf{(1)} $\sin^2\theta=1-9/25=16/25\Rightarrow\sin\theta=\pm4/5$.
\textbf{(2)} $\tan\theta=\pm(4/5)/(3/5)=\pm4/3$.
\textbf{(3)} $\sin\theta=4/5$~: Q1 ($\theta\approx0{,}93$~rad)~;
$\sin\theta=-4/5$~: Q4 ($\theta\approx5{,}36$~rad).
\textbf{(4)} $\sin\theta<0$ et $\cos\theta=3/5$~: Q4~; $\theta=2\pi-\arccos(3/5)\approx5{,}36$~rad.

%────────────────────────────────────────────────────────────────
\psolutiontitle{8}{7}{Sch\'ema du cercle et arc}

\textbf{(1)} $1$~s~: $60/60=1$~tour/s.
\textbf{(2)} $\omega=2\pi\times1=2\pi$~rad/s.
\textbf{(3)} Distance~: $\ell=r\omega t=0{,}30\times2\pi\times1=0{,}6\pi\approx1{,}88$~m.
\textbf{(4)} En $2$~s~: $\omega t=4\pi$~rad.

%────────────────────────────────────────────────────────────────
\psolutiontitle{8}{8}{\'Equations biqu\-adratiques}

\textbf{(1)} $\cos^2x=3/4\Rightarrow\cos x=\pm\frac{\sqrt{3}}{2}$~:
$x=\frac{\pi}{6},\frac{5\pi}{6},\frac{7\pi}{6},\frac{11\pi}{6}$.

\textbf{(2)} $u=\sin x$~: $2u^2+u-1=(2u-1)(u+1)=0\Rightarrow u=\frac{1}{2}$ ou $u=-1$.
$\sin x=\frac{1}{2}$~: $x=\frac{\pi}{6},\frac{5\pi}{6}$~; $\sin x=-1$~: $x=\frac{3\pi}{2}$.

\textbf{(3)} $u=\cos x$~: $u^2+3u+2=(u+1)(u+2)=0\Rightarrow u=-1$ ou $u=-2$.
$\cos x=-2$~: impossible. $\cos x=-1$~: $x=\pi$.

\textbf{(4)} $\sin^2x\in[0,1]$~; $\sin^2x=2$ est impossible.\hfill$\square$

%────────────────────────────────────────────────────────────────
\psolutiontitle{8}{9}{Trigonom\'etrie et g\'eom\'etrie}

$A(0,0)$, $B(4,0)$, $C(1,3)$.
\textbf{(1)} $AB=4$~; $BC=\sqrt{9+9}=3\sqrt{2}$~; $CA=\sqrt{1+9}=\sqrt{10}$.

\textbf{(2)} $\vv{AB}=(4,0)$, $\vv{AC}=(1,3)$.
$\cos\hat{A}=\frac{4}{4\sqrt{10}}=\frac{1}{\sqrt{10}}$~; $\hat{A}=\arccos(1/\sqrt{10})\approx71{,}6^{\circ}$.
$\vv{BA}=(-4,0)$, $\vv{BC}=(-3,3)$.
$\cos\hat{B}=\frac{12}{4\cdot3\sqrt{2}}=\frac{12}{12\sqrt{2}}=\frac{1}{\sqrt{2}}$~; $\hat{B}=45^{\circ}$.
$\hat{C}=180^{\circ}-71{,}6^{\circ}-45^{\circ}=63{,}4^{\circ}$.

\textbf{(3)} $71{,}6^{\circ}+45^{\circ}+63{,}4^{\circ}=180^{\circ}$.\checkmark

\textbf{(4)} Loi des cosinus~: $BC^2=AB^2+CA^2-2\cdot AB\cdot CA\cdot\cos\hat{A}$~?
$18=16+10-2\cdot4\cdot\sqrt{10}\cdot\frac{1}{\sqrt{10}}=26-8=18$.\checkmark

%────────────────────────────────────────────────────────────────
\psolutiontitle{8}{10}{\'Equation trigonom\'etrique appliqu\'ee}

$x(\theta)=3\cos\theta+4\sin\theta$.
\textbf{(1)} $\sqrt{9+16}=5$.
\textbf{(2)} $|x(\theta)|=|3\cos\theta+4\sin\theta|\leq|3\cos\theta|+|4\sin\theta|\leq3+4=7$... mieux~:
par Cauchy-Schwarz, $|3\cos\theta+4\sin\theta|\leq\sqrt{9+16}\sqrt{\cos^2\theta+\sin^2\theta}=5$.\checkmark
\textbf{(3)} $x(\frac{\pi}{3})=3/2+4\cdot\frac{\sqrt{3}}{2}=3/2+2\sqrt{3}$.
\textbf{(4)} $3\cos\theta=-4\sin\theta\Rightarrow\tan\theta=-3/4$.
$\theta=\arctan(-3/4)$~; sur $[0,2\pi]$~: $\theta=\pi-\arctan(3/4)\approx2{,}50$ et $\theta=2\pi-\arctan(3/4)\approx5{,}64$.

%────────────────────────────────────────────────────────────────
\psolutiontitle{8}{11}{Trigonom\'etrie et repr\'esentation}

\textbf{(1)} Tableau pour $\theta_k=k\frac{\pi}{6}$~: voir cours.
\textbf{(2)} Courbes sur $[0,2\pi]$ (voir figure du cours).
\textbf{(3)} $\cos$~: p\'eriode $2\pi$, max $1$ en $0$, min $-1$ en $\pi$, z\'eros en $\frac{\pi}{2}$ et $\frac{3\pi}{2}$.
$\sin$~: p\'eriode $2\pi$, max $1$ en $\frac{\pi}{2}$, min $-1$ en $\frac{3\pi}{2}$, z\'eros en $0,\pi,2\pi$.
\textbf{(4)}
\begin{lstlisting}[language=Python]
import numpy as np, matplotlib.pyplot as plt
t = np.linspace(0, 2*np.pi, 300)
plt.plot(t, np.cos(t), label='cos'); plt.plot(t, np.sin(t), label='sin')
plt.legend(); plt.grid(True); plt.show()
\end{lstlisting}

%────────────────────────────────────────────────────────────────
\psolutiontitle{8}{12}{Longueur d'arc et secteur angulaire}

\textbf{(1)} $A=\frac{1}{2}\cdot64\cdot\frac{\pi}{3}=\frac{32\pi}{3}$~cm$^2$.
\textbf{(2)} $A=\frac{1}{2}r^2\theta\Rightarrow\theta=\frac{2A}{r^2}=\frac{24\pi}{36}=\frac{2\pi}{3}$~rad $=120^{\circ}$.
\textbf{(3)} $\ell=5\cdot\frac{2\pi}{3}=\frac{10\pi}{3}$~cm.
\textbf{(4)} $\theta=\ell/r=15\pi/10=\frac{3\pi}{2}$~rad~; $A=\frac{1}{2}\cdot100\cdot\frac{3\pi}{2}=75\pi$~cm$^2$.

%────────────────────────────────────────────────────────────────
\psolutiontitle{8}{13}{\logiq\ Logique et identit\'es trigonom\'etriques}

\textbf{(1)} $M(\cos\theta,\sin\theta)$ est sur le cercle unit\'e~: $\Omega M=1$, donc
$\cos^2\theta+\sin^2\theta=1$ par d\'efinition de la distance.\hfill$\square$
\textbf{(2)} De $\cos^2\theta+\sin^2\theta=1\geq0$ et $\cos^2\theta\geq0$~: $\sin^2\theta\leq1$~;
prendre la racine~: $|\sin\theta|\leq1$. Idem pour $|\cos\theta|\leq1$.\hfill$\square$
\textbf{(3)} $\cos^4\theta-\sin^4\theta=(\cos^2\theta)^2-(\sin^2\theta)^2=(\cos^2\theta-\sin^2\theta)(\cos^2\theta+\sin^2\theta)=\cos^2\theta-\sin^2\theta$.
Raisonnement~: \textbf{calcul alg\'ebrique direct}.\hfill$\square$

%────────────────────────────────────────────────────────────────
\psolutiontitle{8}{14}{\logiq\ R\'esolution logique d'\'equations compos\'ees}

\textbf{(1)} $u=\sin x$~: $2u^2-3u+1=(2u-1)(u-1)=0$. Loi du produit nul~: $u=\frac{1}{2}$ ou $u=1$.
Solutions~: $x=\frac{\pi}{6},\frac{5\pi}{6},\frac{\pi}{2}$.

\textbf{(2)} $\cos^2x-\sin^2x=(\cos x-\sin x)(\cos x+\sin x)=0$~:
$\cos x=\sin x\Rightarrow x=\frac{\pi}{4},\frac{5\pi}{4}$~; $\cos x=-\sin x\Rightarrow x=\frac{3\pi}{4},\frac{7\pi}{4}$.

\textbf{(3)} $(2\cos x-1)(\sin x+1)=0$~: loi du produit nul~;
$\cos x=\frac{1}{2}$~: $x=\frac{\pi}{3},\frac{5\pi}{3}$~; $\sin x=-1$~: $x=\frac{3\pi}{2}$.

\textbf{(4)} Raisonnements~: (1) substitution + discriminant~; (2) factorisation~; (3) produit nul.

%────────────────────────────────────────────────────────────────
\psolutiontitle{8}{15}{\algo\ Approximation de $\pi$ par la m\'ethode d'Archim\`ede \computer}

\begin{lstlisting}[language=Python]
import math

print(f"pi = {math.pi:.10f}")
for n in [6, 12, 24, 48, 96]:
    Pn = 2 * n * math.sin(math.pi / n)   # inscrit
    Pn_circ = 2 * n * math.tan(math.pi / n)  # circonscrit
    print(f"n={n:3d}: {Pn:.8f} < pi < {Pn_circ:.8f}")
\end{lstlisting}
Pour $n=96$~: $3{,}14103\ldots<\pi<3{,}14271\ldots$~: l'encadrement d'Archim\`ede.

%────────────────────────────────────────────────────────────────
\psolutiontitle{8}{16}{\algo\ D\'eterminer le quadrant \computer}

\begin{lstlisting}[language=Python]
import math

def quadrant(theta):
    t = theta % (2*math.pi)
    c, s = math.cos(t), math.sin(t)
    if c > 0 and s > 0: q = 1
    elif c < 0 and s > 0: q = 2
    elif c < 0 and s < 0: q = 3
    else: q = 4
    print(f"theta={t:.4f}: Q{q}, cos={c:.4f}, sin={s:.4f}")

quadrant(math.pi/4); quadrant(2*math.pi/3); quadrant(5*math.pi/4)
\end{lstlisting}

%────────────────────────────────────────────────────────────────
\subsection*{D\'efis}
\addcontentsline{toc}{subsection}{D\'efis}

%────────────────────────────────────────────────────────────────
\noindent\phantomsection
{\color{BO}\bfseries\small D\'efi~1~\textendash\ \textit{Applications g\'eom\'etriques du cercle}}
\par\smallskip

\textbf{(1)} Q3~: $\theta\in]\pi,\frac{3\pi}{2}[$~; Q4~: $\theta\in]\frac{3\pi}{2},2\pi[$.
\textbf{(2)} $\sin\theta=\cos\theta\Rightarrow\theta=\frac{\pi}{4}$ ou $\theta=\frac{5\pi}{4}$.
\textbf{(3)} $\cos\theta>\sin\theta\Rightarrow x\in[0,\frac{\pi}{4}[\cup]\frac{5\pi}{4},2\pi]$.
\textbf{(4)} $\cos\theta=-\frac{\sqrt{3}}{2}$~: $\theta=\frac{5\pi}{6}$ ou $\theta=\frac{7\pi}{6}$.

%────────────────────────────────────────────────────────────────
\noindent\phantomsection
{\color{BO}\bfseries\small D\'efi~2~\textendash\ \textit{Preuve de la relation de Pythagore trigonom\'etrique}}
\par\smallskip

\textbf{(1)} $\cos\theta=b/c$, $\sin\theta=a/c$.
\textbf{(2)} $\cos^2\theta+\sin^2\theta=b^2/c^2+a^2/c^2=(a^2+b^2)/c^2=c^2/c^2=1$ (par Pythagore).\hfill$\square$
\textbf{(3)} Pour tout $\theta$~: $M(\cos\theta,\sin\theta)$ est \`a distance $1$ de l'origine
(d\'efinition du cercle unit\'e)~: $\cos^2\theta+\sin^2\theta=1$ par d\'efinition m\^eme.\hfill$\square$

%────────────────────────────────────────────────────────────────
\noindent\phantomsection
{\color{BO}\bfseries\small D\'efi~3~\textendash\ \textit{\'Equation $a\cos x+b\sin x=c$}}
\par\smallskip

$\cos x+\sin x=1$.
\textbf{(1)} $(\cos x+\sin x)^2=\cos^2x+2\sin x\cos x+\sin^2x=1+2\sin x\cos x=1$.
\textbf{(2)} $2\sin x\cos x=0$.
\textbf{(3)} $\sin x=0$ ou $\cos x=0$~:
$x=0,\pi,2\pi$ ou $x=\frac{\pi}{2},\frac{3\pi}{2}$.
\textbf{(4)} V\'erification~:
$x=0$~: $1+0=1$.\checkmark\quad
$x=\frac{\pi}{2}$~: $0+1=1$.\checkmark\quad
$x=\pi$~: $-1+0=-1\ne1$~: parasite.\quad
$x=\frac{3\pi}{2}$~: $0+(-1)=-1\ne1$~: parasite.
Solutions~: $x=0$ et $x=\frac{\pi}{2}$.

%────────────────────────────────────────────────────────────────
\noindent\phantomsection
{\color{BO}\bfseries\small D\'efi~4~\textendash\ \textit{Trigonom\'etrie et g\'eom\'etrie~: hexagone r\'egulier}}
\par\smallskip

\textbf{(1)} Angles $0,\frac{\pi}{3},\frac{2\pi}{3},\pi,\frac{4\pi}{3},\frac{5\pi}{3}$~; coordonn\'ees~:
$(1,0)$~; $(\frac{1}{2},\frac{\sqrt{3}}{2})$~; $(-\frac{1}{2},\frac{\sqrt{3}}{2})$~;
$(-1,0)$~; $(-\frac{1}{2},-\frac{\sqrt{3}}{2})$~; $(\frac{1}{2},-\frac{\sqrt{3}}{2})$.

\textbf{(2)} Longueur d'un c\^ot\'e entre angles $0$ et $\frac{\pi}{3}$~:
$MN^2=2-2\cos(\frac{\pi}{3})=2-1=1$~; $MN=1=r$.\hfill$\square$

\textbf{(3)} P\'erim\`etre~: $6r=6$. Aire~: $6$ triangles \'equilat\'eraux de c\^ot\'e $1$~:
$6\times\frac{\sqrt{3}}{4}=\frac{3\sqrt{3}}{2}$.

\textbf{(4)} P\'erim\`etre $=6<2\pi\approx6{,}28$~:
encadrement~: $6<2\pi$, donc $\pi>3$ (borne inf\'erieure d'Archim\`ede).

%────────────────────────────────────────────────────────────────
\noindent\phantomsection
{\color{BO}\bfseries\small D\'efi~5~\textendash\ \textit{Identit\'es et applications \`a la distance}}
\par\smallskip

\textbf{(1)} $(\cos\alpha-\cos\beta)^2+(\sin\alpha-\sin\beta)^2
=2-2(\cos\alpha\cos\beta+\sin\alpha\sin\beta)$~: voir exo~35.\hfill$\square$
\textbf{(2)} $MN=\sqrt{2-2c}$ avec $c=\cos\alpha\cos\beta+\sin\alpha\sin\beta$.
\textbf{(3)} $\alpha=0$, $\beta=\frac{\pi}{3}$~: $c=\cos(\frac{\pi}{3})=\frac{1}{2}$~; $MN=\sqrt{2-1}=1$.\checkmark
\textbf{(4)} $c$ max en $\alpha=\beta$~: $c=\cos^2\alpha+\sin^2\alpha=1$~; $MN=0$ (m\^eme point).

%────────────────────────────────────────────────────────────────
\noindent\phantomsection
{\color{BO}\bfseries\small D\'efi~6~\textendash\ \textit{Trigonom\'etrie et algorithme \computer}}
\par\smallskip

\begin{lstlisting}[language=Python]
import math, numpy as np, matplotlib.pyplot as plt

# (1) Tableau des valeurs
angles = [k*math.pi/6 for k in range(9)]
for t in angles:
    c, s = math.cos(t), math.sin(t)
    print(f"theta={t:.4f}: cos={c:.4f}, sin={s:.4f}, c^2+s^2={c**2+s**2:.10f}")

# (2) Solutions de cos x = 0.5 dans [0, 2pi]
a = 0.5
prev = math.cos(0) - a
for i in range(1, 629):
    x = i * 0.01
    curr = math.cos(x) - a
    if prev * curr < 0:
        print(f"Racine ~ x={x:.2f}")
    prev = curr
# Exact: pi/3 ~ 1.047, 5pi/3 ~ 5.236

# (3) Trace du cercle avec 12 angles
theta_all = np.linspace(0, 2*np.pi, 300)
plt.figure(figsize=(6,6))
plt.plot(np.cos(theta_all), np.sin(theta_all), 'b-')
for k in range(12):
    t = k*np.pi/6
    plt.plot(np.cos(t), np.sin(t), 'ro', ms=6)
plt.axis('equal'); plt.grid(True); plt.show()
\end{lstlisting}

%────────────────────────────────────────────────────────────────
\noindent\phantomsection
{\color{BO}\bfseries\small D\'efi~7~\textendash\ \textit{\logiq\ Propri\'et\'es logiques de la p\'eriodicit\'e}}
\par\smallskip

\textbf{(1)} Sur le cercle, avancer de $2\pi$ revient au m\^eme point~: $\cos(\theta+2\pi)=\cos\theta$.\hfill$\square$
\textbf{(2)} $T$ est p\'eriode de $f\Leftrightarrow\forall x,\;f(x+T)=f(x)$.
La p\'eriode minimale est unique (la \emph{p\'eriode fondamentale})~; $2\pi$ est fondamentale pour $\cos$ et $\sin$.
\textbf{(3)} Si $f(x+T)=f(x)$, alors $f(x+2T)=f((x+T)+T)=f(x+T)=f(x)$~: $2T$ est aussi p\'eriode.\hfill$\square$
\textbf{(4)} Par r\'ecurrence~: $f(x+nT)=f(x+(n-1)T+T)=f(x+(n-1)T)=\ldots=f(x)$.\hfill$\square$

%────────────────────────────────────────────────────────────────
\noindent\phantomsection
{\color{BO}\bfseries\small D\'efi~8~\textendash\ \textit{\algo\ Approcher $\cos$ par un polyn\^ome \computer}}
\par\smallskip

\textbf{(1)} En $x=0$~: $1-0+0=1=\cos(0)$.\checkmark\quad Coh\'erent car le cosinus vaut $1$ en $0$.

\textbf{(2--3)}
\begin{lstlisting}[language=Python]
import math

def cos_approx(x):
    return 1 - x**2/2 + x**4/24

angles = [0, math.pi/6, math.pi/4, math.pi/3, math.pi/2]
for t in angles:
    approx = cos_approx(t)
    exact  = math.cos(t)
    erreur = abs(exact - approx)
    print(f"x={t:.4f}: approx={approx:.6f}, exact={exact:.6f}, erreur={erreur:.2e}")
\end{lstlisting}

\textbf{(4)} En ajoutant $-x^6/720$~:
\begin{lstlisting}[language=Python]
def cos_approx2(x):
    return 1 - x**2/2 + x**4/24 - x**6/720

for t in angles:
    e2 = abs(math.cos(t) - cos_approx2(t))
    e1 = abs(math.cos(t) - cos_approx(t))
    print(f"x={t:.4f}: erreur3={e1:.2e}, erreur4={e2:.2e}")
# L'erreur diminue, surtout pour x proche de pi/2
\end{lstlisting}
L'approximation s'am\'eliore nettement pour $x$ proche de $\frac{\pi}{2}$.\hfill$\square$

\needspace{3\baselineskip}
